

\begin{register}{H}{GPIO\_OUT\_REG}{0x0004}\label{regdesc:GPIOOUTREG}
\regfield{GPIO\_OUT\_DATA\_ORIG}{32}{0}{{0x000000}}%
\reglabel{Reset}\regnewline%
\begin{regdesc}\begin{reglist}
\label{fielddesc:GPIOOUTDATAORIG}\item [GPIO\_OUT\_DATA\_ORIG] Configures the output value of GPIO0 \verb+~+ 30 output in simple GPIO output mode.\\
0: Low level\\
1: High level\\
The value of bit0 \verb+~+ bit30 correspond to the output value of GPIO0 \verb+~+ GPIO30 respectively. Bit31 is invalid.\\
(R/W/SC/WTC)
\end{reglist}\end{regdesc}
\end{register}


\begin{register}{H}{GPIO\_OUT\_W1TS\_REG}{0x0008}\label{regdesc:GPIOOUTW1TSREG}
\regfield{GPIO\_OUT\_W1TS}{32}{0}{{0x000000}}%
\reglabel{Reset}\regnewline%
\begin{regdesc}\begin{reglist}
\label{fielddesc:GPIOOUTW1TS}\item [GPIO\_OUT\_W1TS] Configures whether or not to set the output register \hyperref[regdesc:GPIOOUTREG]{GPIO\_OUT\_REG} of GPIO0 \verb+~+ GPIO30.\\
0: Not set\\
1: The corresponding bit in \hyperref[regdesc:GPIOOUTREG]{GPIO\_OUT\_REG} will be set to 1\\
Bit0 \verb+~+ bit30 are corresponding to GPIO0 \verb+~+ GPIO30. Bit31 is invalid.
Recommended operation: use this register to set \hyperref[regdesc:GPIOOUTREG]{GPIO\_OUT\_REG}. \\
(WT)
\end{reglist}\end{regdesc}
\end{register}


\begin{register}{H}{GPIO\_OUT\_W1TC\_REG}{0x000C}\label{regdesc:GPIOOUTW1TCREG}
\regfield{GPIO\_OUT\_W1TC}{32}{0}{{0x000000}}%
\reglabel{Reset}\regnewline%
\begin{regdesc}\begin{reglist}
\label{fielddesc:GPIOOUTW1TC}\item [GPIO\_OUT\_W1TC] Configures whether or not to clear the output register \hyperref[regdesc:GPIOOUTREG]{GPIO\_OUT\_REG} of GPIO0 \verb+~+ GPIO30 output.\\
0: Not clear\\
1: The corresponding bit in \hyperref[regdesc:GPIOOUTREG]{GPIO\_OUT\_REG} will be cleared.\\
Bit0 \verb+~+ bit30 are corresponding to GPIO0 \verb+~+ GPIO30. Bit31 is invalid.
Recommended operation: use this register to clear \hyperref[regdesc:GPIOOUTREG]{GPIO\_OUT\_REG}. \\
(WT)
\end{reglist}\end{regdesc}
\end{register}




\begin{register}{H}{GPIO\_ENABLE\_REG}{0x0020}\label{regdesc:GPIOENABLEREG}
\regfield{GPIO\_ENABLE\_DATA}{32}{0}{{0x000000}}%
\reglabel{Reset}\regnewline%
\begin{regdesc}\begin{reglist}
\label{fielddesc:GPIOENABLEDATA}\item [GPIO\_ENABLE\_DATA] Configures whether or not to enable the output of GPIO0 \verb+~+ GPIO30.\\
0: Not enable\\
1: Enable\\
Bit0 \verb+~+ bit30 are corresponding to GPIO0 \verb+~+ GPIO30. Bit31 is invalid.\\
(R/W/WTC)
\end{reglist}\end{regdesc}
\end{register}


\begin{register}{H}{GPIO\_ENABLE\_W1TS\_REG}{0x0024}\label{regdesc:GPIOENABLEW1TSREG}
\regfield{GPIO\_ENABLE\_W1TS}{32}{0}{{0x000000}}%
\reglabel{Reset}\regnewline%
\begin{regdesc}\begin{reglist}
\label{fielddesc:GPIOENABLEW1TS}\item [GPIO\_ENABLE\_W1TS] Configures whether or not to set the output enable register \hyperref[regdesc:GPIOENABLEREG]{GPIO\_ENABLE\_REG} of GPIO0 \verb+~+ GPIO30.\\
0: Not set\\
1: The corresponding bit in \hyperref[regdesc:GPIOENABLEREG]{GPIO\_ENABLE\_REG} will be set to 1\\
Bit0 \verb+~+ bit30 are corresponding to GPIO0 \verb+~+ GPIO30. Bit31 is invalid. Recommended operation: use this register to set \hyperref[regdesc:GPIOENABLEREG]{GPIO\_ENABLE\_REG}.\\
(WT)
\end{reglist}\end{regdesc}
\end{register}


\begin{register}{H}{GPIO\_ENABLE\_W1TC\_REG}{0x0028}\label{regdesc:GPIOENABLEW1TCREG}
\regfield{GPIO\_ENABLE\_W1TC}{32}{0}{{0x000000}}%
\reglabel{Reset}\regnewline%
\begin{regdesc}\begin{reglist}
\label{fielddesc:GPIOENABLEW1TC}\item [GPIO\_ENABLE\_W1TC] Configures whether or not to clear the output enable register \hyperref[regdesc:GPIOENABLEREG]{GPIO\_ENABLE\_REG} of GPIO0 \verb+~+ GPIO30. \\
0: Not clear\\
1: The corresponding bit in \hyperref[regdesc:GPIOENABLEREG]{GPIO\_ENABLE\_REG} will be cleared\\
Bit0 \verb+~+ bit30 are corresponding to GPIO0 \verb+~+ 30. Bit31 is invalid. Recommended operation: use this register to clear \hyperref[regdesc:GPIOENABLEREG]{GPIO\_ENABLE\_REG}.\\
(WT)
\end{reglist}\end{regdesc}
\end{register}


\begin{register}{H}{GPIO\_STRAP\_REG}{0x0038}\label{regdesc:GPIOSTRAPREG}
\regfield{(reserved)}{16}{16}{0000000000000000}%
\regfield{GPIO\_STRAPPING}{16}{0}{{0x00}}%
\reglabel{Reset}\regnewline%
\begin{regdesc}\begin{reglist}
\label{fielddesc:GPIOSTRAPPING}\item [GPIO\_STRAPPING] Represents the values of GPIO strapping pins. 
\begin{itemize}
    \item bit0 \verb+~+ bit1: invalid
    \item bit2: GPIO8
    \item bit3: GPIO9
    \item bit4: GPIO15
    \item bit5: MTMS
    \item bit6: MTDI
    \item bit7 \verb+~+ bit15: invalid
\end{itemize}
(RO)
\end{reglist}\end{regdesc}
\end{register}


\begin{register}{H}{GPIO\_IN\_REG}{0x003C}\label{regdesc:GPIOINREG}
\regfield{GPIO\_IN\_DATA\_NEXT}{32}{0}{{0x000000}}%
\reglabel{Reset}\regnewline%
\begin{regdesc}\begin{reglist}
\label{fielddesc:GPIOINDATANEXT}\item [GPIO\_IN\_DATA\_NEXT] Represents the input value of GPIO0 \verb+~+ GPIO30. Each bit represents a pin input value:\\
0: Low level\\
1: High level\\
Bit0 \verb+~+ bit30 are corresponding to GPIO0 \verb+~+ GPIO30. Bit31 is invalid.\\
(RO)
\end{reglist}\end{regdesc}
\end{register}


\begin{register}{H}{GPIO\_STATUS\_REG}{0x0044}\label{regdesc:GPIOSTATUSREG}
\regfield{GPIO\_STATUS\_INTERRUPT}{32}{0}{{0x000000}}%
\reglabel{Reset}\regnewline%
\begin{regdesc}\begin{reglist}
\label{fielddesc:GPIOSTATUSINTERRUPT}\item [GPIO\_STATUS\_INTERRUPT] The interrupt status of GPIO0 \verb+~+ GPIO30, can be configured by the software. 
\begin{itemize}
    \item Bit0 \verb+~+ bit30 are corresponding to GPIO0 \verb+~+ GPIO30. Bit31 is invalid. 
    \item Each bit represents the status of its corresponding GPIO:
    \begin{itemize}
        \item 0: Represents the GPIO does not generate the interrupt configured by \hyperref[fielddesc:GPIOPINNINTTYPE]{GPIO\_PIN\regindex{n}\_INT\_TYPE}, or this bit is configured to 0 by the software.
        \item 1: Represents the GPIO generates the interrupt configured by \hyperref[fielddesc:GPIOPINNINTTYPE]{GPIO\_PIN\regindex{n}\_INT\_TYPE}, or this bit is configured to 1 by the software.
    \end{itemize}
\end{itemize}
(R/W/WTC)
\end{reglist}\end{regdesc}
\end{register}


\begin{register}{H}{GPIO\_STATUS\_W1TS\_REG}{0x0048}\label{regdesc:GPIOSTATUSW1TSREG}
\regfield{GPIO\_STATUS\_W1TS}{32}{0}{{0x000000}}%
\reglabel{Reset}\regnewline%
\begin{regdesc}\begin{reglist}
\label{fielddesc:GPIOSTATUSW1TS}\item [GPIO\_STATUS\_W1TS] Configures whether or not to set the interrupt status register \hyperref[fielddesc:GPIOSTATUSINTERRUPT]{GPIO\_STATUS\_INTERRUPT} of GPIO0 \verb+~+ GPIO30. 
\begin{itemize}
    \item Bit0 \verb+~+ bit30 are corresponding to GPIO0 \verb+~+ GPIO30. Bit31 is invalid.
    \item If the value 1 is written to a bit here, the corresponding bit in \hyperref[fielddesc:GPIOSTATUSINTERRUPT]{GPIO\_STATUS\_INTERRUPT} will be set to 1. \item Recommended operation: use this register to set \hyperref[fielddesc:GPIOSTATUSINTERRUPT]{GPIO\_STATUS\_INTERRUPT}. 
\end{itemize}
(WT)
\end{reglist}\end{regdesc}
\end{register}


\begin{register}{H}{GPIO\_STATUS\_W1TC\_REG}{0x004C}\label{regdesc:GPIOSTATUSW1TCREG}
\regfield{GPIO\_STATUS\_W1TC}{32}{0}{{0x000000}}%
\reglabel{Reset}\regnewline%
\begin{regdesc}\begin{reglist}
\label{fielddesc:GPIOSTATUSW1TC}\item [GPIO\_STATUS\_W1TC] Configures whether or not to clear the interrupt status register \hyperref[fielddesc:GPIOSTATUSINTERRUPT]{GPIO\_STATUS\_INTERRUPT} of GPIO0 \verb+~+ GPIO30. 
\begin{itemize}
    \item Bit0 \verb+~+ bit30 are corresponding to GPIO0 \verb+~+ GPIO30. Bit31 is invalid.
    \item If the value 1 is written to a bit here, the corresponding bit in \hyperref[fielddesc:GPIOSTATUSINTERRUPT]{GPIO\_STATUS\_INTERRUPT} will be cleared. \item Recommended operation: use this register to clear \hyperref[fielddesc:GPIOSTATUSINTERRUPT]{GPIO\_STATUS\_INTERRUPT}. \end{itemize}(WT)
\end{reglist}\end{regdesc}
\end{register}


\begin{register}{H}{GPIO\_PCPU\_INT\_REG}{0x005C}\label{regdesc:GPIOPCPUINTREG}
\regfield{GPIO\_PROCPU\_INT}{32}{0}{{0x000000}}%
\reglabel{Reset}\regnewline%
\begin{regdesc}\begin{reglist}
\label{fielddesc:GPIOPROCPUINT}\item [GPIO\_PROCPU\_INT] Represents the CPU interrupt status of GPIO0 \verb+~+ GPIO30. Each bit represents:\\
0: Represents CPU interrupt is not enabled, or the GPIO does not generate the interrupt configured by \hyperref[fielddesc:GPIOPINNINTTYPE]{GPIO\_PIN\regindex{n}\_INT\_TYPE}.\\
1: Represents the GPIO generates an interrupt configured by \hyperref[fielddesc:GPIOPINNINTTYPE]{GPIO\_PIN\regindex{n}\_INT\_TYPE} after the CPU interrupt is enabled.\\
Bit0 \verb+~+ bit30 are corresponding to GPIO0 \verb+~+ GPIO30. Bit31 is invalid. This interrupt status is corresponding to the bit in \hyperref[regdesc:GPIOSTATUSREG]{GPIO\_STATUS\_REG} when assert (high) enable signal (bit13 of \hyperref[regdesc:GPIOPINNREG]{GPIO\_PIN\regindex{n}\_REG}). \\
(RO)
\end{reglist}\end{regdesc}
\end{register}








\begin{register}{H}{GPIO\_PIN\regindex{n}\_REG (\regindex{n}: 0-30)}{0x0074+4*\regindex{n}}\label{regdesc:GPIOPINNREG}
\regfield{(reserved)}{14}{18}{00000000000000}%
\regfield{GPIO\_PIN\regindex{n}\_INT\_ENA}{5}{13}{{0x0}}%
\regfield{(reserved)}{2}{11}{{0x0}}%
\regfield{GPIO\_PIN\regindex{n}\_WAKEUP\_ENABLE}{1}{10}{{0}}%
\regfield{GPIO\_PIN\regindex{n}\_INT\_TYPE}{3}{7}{{0x0}}%
\regfield{(reserved)}{2}{5}{00}%
\regfield{GPIO\_PIN\regindex{n}\_SYNC1\_BYPASS}{2}{3}{{0x0}}%
\regfield{GPIO\_PIN\regindex{n}\_PAD\_DRIVER}{1}{2}{{0}}%
\regfield{GPIO\_PIN\regindex{n}\_SYNC2\_BYPASS}{2}{0}{{0x0}}%
\reglabel{Reset}\regnewline%
\begin{regdesc}\begin{reglist}
\label{fielddesc:GPIOPINNSYNC2BYPASS}\item [GPIO\_PIN\regindex{n}\_SYNC2\_BYPASS] Configures whether or not to synchronize GPIO input data on either edge of IO MUX operating clock for the second-level synchronization.\\
0: Not synchronize\\
1: Synchronize on falling edge\\
2: Synchronize on rising edge\\
3: Synchronize on rising edge\\
(R/W)
\label{fielddesc:GPIOPINNPADDRIVER}\item [GPIO\_PIN\regindex{n}\_PAD\_DRIVER] Configures to select pin drive mode. \\
0: Normal output\\
1: Open drain output \\(R/W)
\label{fielddesc:GPIOPINNSYNC1BYPASS}\item [GPIO\_PIN\regindex{n}\_SYNC1\_BYPASS] Configures whether or not to synchronize GPIO input data on either edge of IO MUX operating clock for the first-level synchronization.\\
0: Not synchronize\\
1: Synchronize on falling edge\\
2: Synchronize on rising edge\\
3: Synchronize on rising edge\\
(R/W)
\label{fielddesc:GPIOPINNINTTYPE}\item [GPIO\_PIN\regindex{n}\_INT\_TYPE] Configures GPIO interrupt type.\\
0: GPIO interrupt disabled\\
1: Rising edge trigger\\
2: Falling edge trigger\\
3: Any edge trigger\\
4: Low level trigger\\
5: High level trigger\\ (R/W)
\label{fielddesc:GPIOPINNWAKEUPENABLE}\item [GPIO\_PIN\regindex{n}\_WAKEUP\_ENABLE] Configures whether or not to enable GPIO wake-up function.\\
0: Disable\\
1: Enable\\
This function only wakes up the CPU from Light-sleep. \\(R/W)
\item [Continued on the next page...]
\end{reglist}\end{regdesc}
\end{register}

\addtocounter{Regfloat}{-1}
\begin{register}{H}{GPIO\_PIN\regindex{n}\_REG (\regindex{n}: 0-30)}{0x0074+4*\regindex{n}}
\begin{regdesc}\begin{reglist}
\item [Continued from the previous page...]
\label{fielddesc:GPIOPINNINTENA}\item [GPIO\_PIN\regindex{n}\_INT\_ENA] Configures whether or not to enable CPU interrupt or CPU non-maskable interrupt.
\begin{itemize}
    \item bit13: Configures whether or not to enable CPU interrupt:\\
    0: Disable\\
    1: Enable\\
    \item bit14: Configures CPU non-maskable interrupt: \\
    0: Disable\\
    1: Enable\\
    \item bit15 \verb+~+ bit17: invalid
    \end{itemize} (R/W)
\end{reglist}\end{regdesc}
\end{register}


\begin{register}{H}{GPIO\_STATUS\_NEXT\_REG}{0x014C}\label{regdesc:GPIOSTATUSNEXTREG}
\regfield{GPIO\_STATUS\_INTERRUPT\_NEXT}{32}{0}{{0x000000}}%
\reglabel{Reset}\regnewline%
\begin{regdesc}\begin{reglist}
\label{fielddesc:GPIOSTATUSINTERRUPTNEXT}\item [GPIO\_STATUS\_INTERRUPT\_NEXT] Represents the interrupt source signal of GPIO0 \verb+~+ GPIO30.\\
Bit0 \verb+~+ bit30 are corresponding to GPIO0 \verb+~+ 30. Bit31 is invalid. Each bit represents:\\
0: The GPIO does not generate the interrupt configured by \hyperref[fielddesc:GPIOPINNINTTYPE]{GPIO\_PIN\regindex{n}\_INT\_TYPE}.\\
1: The GPIO generates an interrupt configured by \hyperref[fielddesc:GPIOPINNINTTYPE]{GPIO\_PIN\regindex{n}\_INT\_TYPE}.\\
The interrupt could be rising edge interrupt, falling edge interrupt, level sensitive interrupt and any edge interrupt.\\  (RO)
\end{reglist}\end{regdesc}
\end{register}


\begin{register}{H}{GPIO\_FUNC\regindex{n}\_IN\_SEL\_CFG\_REG (\regindex{n}: 0-127)}{0x0154+4*\regindex{n}}\label{regdesc:GPIOFUNCNINSELCFGREG}
\regfield{(reserved)}{24}{8}{000000000000000000000000}%
\regfield{GPIO\_SIG\regindex{n}\_IN\_SEL}{1}{7}{{0}}%
\regfield{GPIO\_FUNC\regindex{n}\_IN\_INV\_SEL}{1}{6}{{0}}%
\regfield{GPIO\_FUNC\regindex{n}\_IN\_SEL}{6}{0}{{0x0}}%
\reglabel{Reset}\regnewline%
\begin{regdesc}\begin{reglist}
\label{fielddesc:GPIOFUNCNINSEL}\item [GPIO\_FUNC\regindex{n}\_IN\_SEL] Configures to select a pin from the 31 GPIO pins to connect the input signal \regindex{n}.\\
0: Select GPIO0\\
1: Select GPIO1\\
......\\
29: Select GPIO29\\
30: Select GPIO30\\
Or\\
0x38: A constantly high input\\
0x3C: A constantly low input\\
(R/W)
\label{fielddesc:GPIOFUNCNININVSEL}\item [GPIO\_FUNC\regindex{n}\_IN\_INV\_SEL] Configures whether or not to invert the input value.\\
0: Not invert\\
1: Invert\\ (R/W)
\label{fielddesc:GPIOSIGNINSEL}\item [GPIO\_SIG\regindex{n}\_IN\_SEL] Configures whether or not to route signals via GPIO matrix.\\
0: Bypass GPIO matrix, i.e., connect signals directly to peripheral configured in IO MUX.\\
1: Route signals via GPIO matrix.\\  (R/W)
\end{reglist}\end{regdesc}
\end{register}


\begin{register}{H}{GPIO\_FUNC\regindex{n}\_OUT\_SEL\_CFG\_REG (\regindex{n}: 0-30)}{0x0554+4*\regindex{n}}\label{regdesc:GPIOFUNCNOUTSELCFGREG}
\regfield{(reserved)}{21}{11}{000000000000000000000}%
\regfield{GPIO\_FUNC0\_OEN\_INV\_SEL}{1}{10}{{0}}%
\regfield{GPIO\_FUNC0\_OEN\_SEL}{1}{9}{{0}}%
\regfield{GPIO\_FUNC0\_OUT\_INV\_SEL}{1}{8}{{0}}%
\regfield{GPIO\_FUNC0\_OUT\_SEL}{8}{0}{{0x80}}%
\reglabel{Reset}\regnewline%
\begin{regdesc}\begin{reglist}
\label{fielddesc:GPIOFUNCNOUTSEL}\item [GPIO\_FUNC\regindex{n}\_OUT\_SEL] Configures to select a signal \regindex{Y} (0 <= \regindex{Y} < 128) from 128 peripheral signals to be output from GPIO\regindex{n}.\\
0: Select signal 0\\
1: Select signal 1\\
......\\
126: Select signal 126\\
127: Select signal 127\\
Or\\
128: Bit \regindex{n} of \hyperref[regdesc:GPIOOUTREG]{GPIO\_OUT\_REG} and \hyperref[regdesc:GPIOENABLEREG]{GPIO\_ENABLE\_REG} are selected as the output value and output enable.

For the detailed signal list, see Table \ref{tab:iomuxgpio-periph-signals-via-gpio-matrix}.

(R/W/SC)
\label{fielddesc:GPIOFUNCNOUTINVSEL}\item [GPIO\_FUNC\regindex{n}\_OUT\_INV\_SEL] Configures whether or not to invert the output value.\\
0: Not invert\\
1: Invert\\ (R/W/SC)
\label{fielddesc:GPIOFUNCNOENSEL}\item [GPIO\_FUNC\regindex{n}\_OEN\_SEL] Configures to select the source of output enable signal.\\
0: Use output enable signal from peripheral.\\ 
1: Force the output enable signal to be sourced from bit \regindex{n} of \hyperref[regdesc:GPIOENABLEREG]{GPIO\_ENABLE\_REG}. \\(R/W)
\label{fielddesc:GPIOFUNCNOENINVSEL}\item [GPIO\_FUNC\regindex{n}\_OEN\_INV\_SEL] Configures whether or not to invert the output enable signal.\\
0: Not invert\\
1: Invert\\ (R/W)
\end{reglist}\end{regdesc}
\end{register}


\begin{register}{H}{GPIO\_CLOCK\_GATE\_REG}{0x062C}\label{regdesc:GPIOCLOCKGATEREG}
\regfield{(reserved)}{31}{1}{0000000000000000000000000000000}%
\regfield{GPIO\_CLK\_EN}{1}{0}{{1}}%
\reglabel{Reset}\regnewline%
\begin{regdesc}\begin{reglist}
\label{fielddesc:GPIOCLKEN}\item [GPIO\_CLK\_EN] Configures whether or not to enable clock gate.\\
0: Not enable\\
1: Enable, the clock is free running. \\
(R/W)
\end{reglist}\end{regdesc}
\end{register}


\begin{register}{H}{GPIO\_DATE\_REG}{0x06FC}\label{regdesc:GPIOREGDATEREG}
\regfield{(reserved)}{4}{28}{0000}%
\regfield{GPIO\_DATE}{28}{0}{{0x1907040}}%
\reglabel{Reset}\regnewline%
\begin{regdesc}\begin{reglist}
\label{fielddesc:GPIOREGDATE}\item [GPIO\_DATE] Version control register. \\
(R/W)
\end{reglist}\end{regdesc}
\end{register}
